\section*{Critical pruning and the local step identity}

For a nonconstant Boolean function $f$ write
\[
 H=\mathsf H(f),\qquad I=\mathsf I(f),\qquad R=H/I.
\]
Fix a relevant coordinate $i$.  Let $f^+,f^-$ be the two physical
children and put
\[
 \bar H_i=\frac{\mathsf H(f^+)+\mathsf H(f^-)}2,
 \quad
 \bar I_i=\frac{\mathsf I(f^+)+\mathsf I(f^-)}2,
 \quad d_i=\mathsf I_i(f),
 \quad \Delta_iH=H-\bar H_i.
\]
The exact restriction identity is
\[
 I=\bar I_i+d_i.                                      \tag{1}
\]

\begin{lemma}[critical pruning identity]\label{lem:pruning}
For every relevant coordinate,
\[
 \boxed{\bar H_i-R\bar I_i=Rd_i-\Delta_iH.}           \tag{2}
\]
If $\bar I_i>0$ and $s_i:=\Delta_iH/d_i$, then
\[
 \frac{\bar H_i}{\bar I_i}-R
 =\frac{d_i}{\bar I_i}(R-s_i).                        \tag{3}
\]
Consequently, if $s_i<R$, at least one nonconstant physical child has
entropy/influence ratio strictly greater than $R$.
\end{lemma}

\begin{proof}
Equation (2) is obtained by subtracting $R$ times (1) from
$H=\bar H_i+\Delta_iH$.  Equation (3) is division by $\bar I_i$.
Finally, $\bar H_i/\bar I_i$ is the influence-weighted average of the ratios
of the nonconstant children (constant children have zero entropy and
influence).  If this average exceeds $R$, one child ratio exceeds $R$.
\end{proof}

\begin{corollary}[necessary local condition for critical sources]
If every nonconstant physical descendant $g$ of $f$ satisfies
$R(g)\le R(f)$, then every relevant coordinate obeys
\[
 \boxed{\Delta_iH\ge R(f)\,\mathsf I_i(f).}            \tag{4}
\]
More generally, if every descendant has ratio at most $R(f)+\eta$, then
\[
 \Delta_iH\ge R(f)d_i-\eta\bar I_i.                  \tag{5}
\]
\end{corollary}

\begin{proof}
Criticality makes each child ratio at most $R$, hence
$\bar H_i\le R\bar I_i$; use (2).  The same argument with $R+\eta$
gives (5).
\end{proof}

\begin{proposition}[fixed-constant minimal counterexample]
Fix $C<\infty$.  If some Boolean function has
$D_C(f):=\mathsf H(f)-C\mathsf I(f)>0$, choose one of minimum dimension after
deleting irrelevant variables.  Then:
\begin{enumerate}
\item every proper nonconstant physical descendant $g$ has $D_C(g)\le0$;
\item for every coordinate $i$,
\[
 \boxed{\Delta_iH-Cd_i\ge D_C(f)>0;}                  \tag{6}
\]
\item $f$ has no disjoint tensor factorization into two nonconstant Boolean
functions.
\end{enumerate}
\end{proposition}

\begin{proof}
The first claim is minimality.  For every split,
\[
 D_C(f)=\frac{D_C(f^+)+D_C(f^-)}2+\Delta_iH-Cd_i.
\]
The averaged child deficit is nonpositive, proving (6).  For disjoint tensor
products, both entropy and influence are additive, so
$D_C(g\otimes h)=D_C(g)+D_C(h)$.  A positive sum has a positive summand in
smaller dimension, contradicting minimality.
\end{proof}

\paragraph{Substitution into the sparse-activation route.}
The complete saved B01/A005 calculation proves that for the A002 activated
embedding $F_m$, every fresh activation control $z$ satisfies
\[
 \frac{\Delta_zH(F_m)}{\mathsf I_z(F_m)}\longrightarrow4,
 \qquad
 R(F_m)\longrightarrow R(g)
\]
for the embedded seed $g$.  Hence if $R(g)>4+\epsilon$, then for all large
$m$ the control has step ratio strictly below $R(F_m)$.  Lemma
\ref{lem:pruning} supplies an actual higher-ratio child.  In this activation
split the other child is constant, so the active child itself has larger
ratio.  Therefore the explicit A002/A005 hardness embedding preserves global
$H/I$ and signed downward $L/I$, but it cannot be a descendant-critical
endpoint above ratio four.

This does not prove a universal constant four: addressing and point defects
approach the same boundary, and B05 shows that rare local logarithms can be
large outside a critical endpoint.  The remaining exact target is to prove a
dimension-independent upper bound on $\min_i\Delta_iH/d_i$ for
descendant-critical, tensor-prime Boolean sources, or to construct a genuine
critical family where this minimum is unbounded.
